\section{Conclusion}
\label{sec:conclusion}

We presented \wcabench, a standardized benchmark for prediction and inference
on World Cube Association competition data. Built on an official export of
$6{,}909{,}454$ results spanning $298{,}551$ competitors and $18{,}708$
competitions, it formalizes five tasks---result prediction, placement
prediction, \DNF{} prediction, human-limit estimation and skill
transfer---under a single interface with a leakage-safe rolling-window
protocol, frozen statistics, four-dimensional stratification, a full
statistical-testing toolbox with effect sizes, and a documented compute and
hardware specification.

Beyond the infrastructure, the paper reports 21 real reference baselines across
six methodological families, each evaluated both for a single seed and as a
mean over three seeds. These include a strong tree-based result model
(\maelog{} $=0.164$, or $0.1055\pm0.0417$ across seeds) that beats a sequence
model by roughly an order of magnitude, a three-way Kendall $\kendall=0.7673$
placement tie among the psychology sheet, Plackett--Luce and a kernel-density
simulation while a graph network reaches only $0.567$, a tree-based \DNF{}
classifier (\aucpr{} $=0.373$) whose edge over the historical rate is \emph{not}
significant and which is beaten on average by a simple Beta--Binomial prior,
event-level human-limit estimates that disagree sharply with domain anchors,
and a monotone collapse in identified transfer pairs from $413$ (correlation)
to $6$ (difference-in-differences). We deliberately frame these as a
reproducible starting point rather than as state-of-the-art claims.

The benchmark's value is the question it makes answerable: under the real
constraints of a rule-governed sport, how do different methodological families
compare? We release the data pipeline, task definitions, baseline
implementations and leaderboard tooling, and we invite the community to improve
the baselines and to extend the tasks. Future work includes full-scale
rolling-window evaluation on all $1.72$M test records across many seeds,
domain-aware sequence and graph baselines that can exploit round-format and
head-to-head structure, genuinely causal designs for \task{5}, and a community
challenge with a public leaderboard. We hope \wcabench{} both serves
speedcubing and offers a reusable template for trustworthy benchmarking in
other rule-governed athletic domains.
